\documentclass[11pt]{article}

\usepackage[margin=1in]{geometry}
\usepackage{amsmath,amssymb,amsthm}
\usepackage{hyperref}
\usepackage{enumitem}
\usepackage{courier}
\usepackage{graphicx}
\usepackage{mathtools}
\usepackage{booktabs}
\usepackage{listings}
\usepackage{xcolor}

\lstdefinestyle{code}{
  basicstyle=\ttfamily\small,
  keywordstyle=\bfseries\ttfamily\small\color{blue!70!black},
  commentstyle=\itshape\color{gray!70!black},
  stringstyle=\color{green!40!black},
  showstringspaces=false,
  columns=fullflexible,
  frame=single,
  framerule=0.2pt,
  rulecolor=\color{black!30},
  breaklines=true
}

\title{PIRTM--Phase Mirror Integration and Meta-Ensembles\\
\large A Structured Development Report}
\author{Multiplicity Foundation -- Phase Mirror HQ}
\date{\today}

\begin{document}
\maketitle

\tableofcontents

\section{Executive Summary}

This report documents the recent developments around integrating the \emph{Prime-Indexed Recursive Tensor Mathematics} (PIRTM) runtime and dialect into the Phase Mirror governance stack, and establishing a clean baseline for the \texttt{meta-ensembles} package within the Phase Mirror HQ monorepo \cite{pirtm-pyproject,pm-phases0-3,pm-gcp-guide,pm-mvp-analysis,hq-pirtm-local}. The work is guided by Phase Mirror's architectural principles: explicit tensions, governance-first constraints, and open-core viability.

Three main themes run through the work:
\begin{enumerate}[label=(\alph*)]
  \item Establishing a \emph{runtime-first} PIRTM integration path for Phase Mirror's Python MVP, with MLIR support staged as an optional, higher-complexity backend.
  \item Using the Phase Mirror HQ monorepo to prototype tight coupling (local PIRTM, meta-ensembles) while explicitly aiming at an external \texttt{pirtm} package as the long-term dependency boundary.
  \item Formalizing these decisions into artifacts: adapter modules, tests, ADRs, docs, and configuration, rather than implicit conventions.
\end{enumerate}

The central tension is always made explicit: we trade off expressive power (MLIR, meta-ensembles wiring) and research velocity against reliability, portability, and governance guarantees for Phase Mirror users \cite{pm-phases0-3,pm-mvp-analysis}. The report concludes with a set of precision questions and recommended next steps that preserve Phase Mirror's L0-style invariants while allowing PIRTM to serve as a mathematical motor for meta-governance and dissonance analysis.

\section{Background and Context}

\subsection{Phase Mirror Architecture}

Phase Mirror is an AI governance tool built around a three-tier hierarchical compute model (L0, L1, L2) and a dissonance-centric methodology \cite{pm-phases0-3,pm-mvp-analysis,pm-gcp-guide}. L0 invariants encode non-negotiable safety and correctness checks (schema alignment, nonce validity, permission configurations, etc.). L1 and L2 provide increasingly expensive reasoning and policy evaluation.

From the open-core perspective, Phase Mirror is structured as a pnpm monorepo with:
\begin{itemize}
  \item A core library package (\texttt{mirror-dissonance}) implementing invariants, stores, adapters and the Oracle.
  \item CLI and MCP server packages for integration with developer workflows and AI coding agents.
  \item Governance documentation (ADRs, success criteria, license analyses).
  \item Infrastructure-as-code for AWS and GCP backends (DynamoDB, Firestore, SSM/Secret Manager, KMS, Cloud Storage) \cite{pm-gcp-guide,pm-open-core-dev}.
\end{itemize}

The Phase Mirror HQ repository extends this into a multi-package Python workspace with:
\begin{itemize}
  \item A Python MVP package (\texttt{phase-mirror-mvp}) mirroring core governance functionality.
  \item A local \texttt{pirtm} package.
  \item A \texttt{meta-ensembles} package that uses PIRTM-like structures to aggregate and ledger ensembles of models or oracles.
\end{itemize}

\subsection{PIRTM: Prime-Indexed Recursive Tensor Mathematics}

The PIRTM repository (\texttt{MultiplicityFoundation/PIRTM}) is a Python package with the following core characteristics \cite{pirtm-pyproject}:
\begin{itemize}
  \item Package name \texttt{pirtm}, version \texttt{0.4.1}, with runtime implemented in pure Python and NumPy.
  \item Focus on contractive recurrence dynamics (prime-indexed recursion over tensor states) and a compile-time MLIR dialect for advanced optimization and compilation.
  \item A \texttt{pyproject.toml} that:
    \begin{itemize}
      \item Uses \texttt{setuptools} as build backend.
      \item Requires Python $\geq 3.10$.
      \item Exposes a CLI entrypoint \texttt{pirtm = pirtm.transpiler.cli:main}.
      \item Declares optional extra dependencies for spectral analysis and development.
      \item Registers subpackages: backend, bindings, channels, core, dialect, integrations, mlir, spectral, transpiler, tools, type\_inference.
    \end{itemize}
\end{itemize}

In the Phase Mirror HQ workspace, a local \texttt{packages/pirtm} mirrors the structure and enables tighter integration experiments \cite{hq-pirtm-local}. This local package is used both by \texttt{phase-mirror-mvp} and \texttt{meta-ensembles}.

\subsection{Meta-Ensembles Package}

The \texttt{meta-ensembles} package is a Python package within the Phase Mirror HQ monorepo that implements:
\begin{itemize}
  \item An ensemble ledger and aggregator for combining multiple models or oracles.
  \item A PIRTM link ensemble implementation that bridges ensemble states with PIRTM-style recurrences.
  \item A CLI for running meta-ensemble operations.
  \item A comprehensive test suite that previously depended on local PIRTM fixtures and import paths.
\end{itemize}

Before the recent work, some tests failed due to unavailable \texttt{pirtm} imports and misaligned fixture paths. The new changes establish a passing baseline.

\section{PIRTM--Phase Mirror Integration}

\subsection{Central Tension}

The major design tension for integration is:
\begin{quote}
\emph{Should Phase Mirror depend directly on both the PIRTM MLIR dialect and the Python runtime, or should it initially integrate only the runtime and treat MLIR as an optional optimization?}
\end{quote}

The final decision is to adopt a \emph{runtime-first, MLIR-optional} strategy:
\begin{itemize}
  \item The Python runtime (NumPy-based) becomes the primary, portable dependency.
  \item The MLIR dialect is staged behind an optional backend flag and separate CI job.
  \item Phase Mirror's governance decisions must never depend solely on MLIR being available.
\end{itemize}

This aligns with existing governance-first choices in Phase Mirror, where L0 invariants are always-on and cloud backends are pluggable via adapters rather than hard requirements \cite{pm-phases0-3,pm-open-core-dev,pm-gcp-guide}.

\subsection{Runtime-Only Bridge}

\subsubsection{PIRTM Repository Setup}

The PIRTM package is configured via the following \texttt{pyproject.toml} fragment \cite{pirtm-pyproject}:

\begin{lstlisting}[style=code,language={[LaTeX]TeX}]
[build-system]
requires = ["setuptools>=61"]
build-backend = "setuptools.build_meta"

[project]
name = "pirtm"
version = "0.4.1"
description = "Prime-Indexed Recursive Tensor Mathematics ..."
requires-python = ">=3.10"
dependencies = ["numpy>=1.24"]

[project.optional-dependencies]
spectral = ["pandas>=1.5"]
dev = ["pytest>=7", "numpy>=1.24"]

[project.scripts]
pirtm = "pirtm.transpiler.cli:main"

[tool.setuptools]
packages = [
  "pirtm",
  "pirtm.backend",
  "pirtm.bindings",
  "pirtm.channels",
  "pirtm.core",
  "pirtm.dialect",
  "pirtm.integrations",
  "pirtm.mlir",
  "pirtm.spectral",
  "pirtm.transpiler",
  "pirtm.tools",
  "pirtm.type_inference",
]
\end{lstlisting}

A minimal setup sequence is:

\begin{lstlisting}[style=code]
git clone https://github.com/MultiplicityFoundation/PIRTM.git
cd PIRTM
python -m venv .venv
source .venv/bin/activate  # Windows: .venv\Scripts\activate
pip install -e ".[dev,spectral]"
pirtm --help
pytest tests/ -v
\end{lstlisting}

\subsubsection{Phase Mirror MVP Adapter}

Within the Phase Mirror HQ monorepo, a runtime-only adapter module \texttt{packages/phase-mirror-mvp/pirtm\_adapter.py} has been introduced. This module:
\begin{itemize}
  \item Imports \texttt{pirtm} (for now from the local \texttt{packages/pirtm}) and exposes a narrow API to Phase Mirror, such as:
  \begin{itemize}
    \item \verb|compile_policy_tensor(...)|
    \item \verb|evaluate_recurrence(...)|
  \end{itemize}
  \item Encapsulates PIRTM-specific details so the rest of Phase Mirror operates over a simple, stable interface.
  \item Is explicitly documented as a runtime-only bridge, with MLIR support to come later.
\end{itemize}

A tiny illustrative code snippet (structurally representative) is:

\begin{lstlisting}[style=code,language=Python]
# packages/phase-mirror-mvp/pirtm_adapter.py

from __future__ import annotations
from typing import Any, Dict

import numpy as np
import pirtm  # local monorepo package for now

def evaluate_recurrence(initial_state: np.ndarray,
                        params: Dict[str, Any],
                        steps: int = 10) -> np.ndarray:
    """
    Evaluate a simple contractive PIRTM recurrence from Phase Mirror.
    """
    state = np.array(initial_state, dtype=float)

    for k in range(steps):
        # Hypothetical PIRTM step; real logic uses pirtm.core APIs
        state = pirtm.core.recurrence.contractive_step(
            state,
            params=params,
            step_index=k,
        )

    return state
\end{lstlisting}

A corresponding unit test validates that:
\begin{itemize}
  \item The adapter can import the local PIRTM package.
  \item A simple recurrence converges or remains bounded under contractive dynamics.
\end{itemize}

\subsubsection{Governance Artifact: ADR-028}

An Architecture Decision Record (\texttt{ADR-028-pirtm-integration.md}) has been added under \texttt{packages/phase-mirror-mvp/governance}, stating that:
\begin{itemize}
  \item PIRTM integration is optional and must fail \emph{degraded}, not \emph{hard};
  \item The Python runtime is the authoritative required layer; MLIR is an enhancement;
  \item Phase Mirror must remain usable without PIRTM for the open-core MVP.
\end{itemize}

This is consistent with earlier ADR-driven decisions about L0 invariants and cloud adapters \cite{pm-phases0-3,pm-open-core-dev,pm-mvp-analysis}.

\subsection{Intermediate Step Toward External \texorpdfstring{\texttt{pirtm}}{pirtm} Dependency}

Within the HQ workspace, the current state is:
\begin{itemize}
  \item \texttt{packages/pirtm} serves as the authoritative PIRTM source for development.
  \item \texttt{phase-mirror-mvp} imports PIRTM via local path hacks in tests, while the production code expects \texttt{import pirtm} like a regular library.
  \item The target state is for \texttt{phase-mirror-mvp} to depend on an external \texttt{pirtm} release (PyPI or VCS URL), with no test-level sys.path manipulation.
\end{itemize}

In terms of dependency configuration, the recommended approach is:
\begin{itemize}
  \item In \texttt{phase-mirror-mvp/requirements.txt} or \texttt{pyproject.toml}, add:
\begin{lstlisting}[style=code]
pirtm @ git+https://github.com/MultiplicityFoundation/PIRTM.git@main
\end{lstlisting}
  \item Later, once a stable version is released:
\begin{lstlisting}[style=code]
pirtm>=0.4.1
\end{lstlisting}
\end{itemize}

The governance constraint is that \emph{PIRTM integration is not a gating factor for Phase Mirror's open-core MVP} \cite{pm-mvp-analysis}. It provides additional analytical power, but core governance value must stand without it.

\subsection{MLIR Backend as Optional Path}

To integrate the PIRTM MLIR dialect safely, a two-job CI structure is suggested:
\begin{enumerate}[label=(\roman*)]
  \item Default job: installs \texttt{pirtm} with no MLIR toolchain, runs all Phase Mirror tests assuming the runtime backend.
  \item MLIR job: installs LLVM/MLIR, enables an environment flag such as \verb|PIRTM_BACKEND=mlir| and runs MLIR-specific tests, potentially marked as allow-failure initially.
\end{enumerate}

\noindent A sketch of such a matrix in GitHub Actions YAML form is:

\begin{lstlisting}[style=code]
jobs:
  test:
    runs-on: ubuntu-latest
    strategy:
      matrix:
        backend: [runtime, mlir]

    steps:
      - uses: actions/checkout@v4

      - uses: actions/setup-python@v5
        with:
          python-version: "3.11"

      - name: Install dependencies
        run: |
          pip install -e ".[dev]"
          if [ "${{ matrix.backend }}" = "mlir" ]; then
            ./scripts/install_mlir_toolchain.sh
          fi
          pip install pirtm>=0.4.1

      - name: Run tests
        env:
          PIRTM_BACKEND: ${{ matrix.backend }}
        run: |
          pytest -q
\end{lstlisting}

The Phase Mirror governance ADR should explicitly state that:
\begin{quote}
  If the MLIR toolchain is missing or fails, Phase Mirror must fall back to the runtime backend, not fail governance checks.
\end{quote}

\section{Meta-Ensembles Baseline Stabilization}

\subsection{Package Overview}

The \texttt{meta-ensembles} package within Phase Mirror HQ has:
\begin{itemize}
  \item \texttt{meta\_ensembles/core/ensemble\_ledger.py}: a ledger for recording model/ensemble judgments, timestamps, and provenance.
  \item \texttt{meta\_ensembles/core/ensemble\_aggregator.py}: aggregation logic over ensembles (e.g., majority vote, weighted averages).
  \item \texttt{meta\_ensembles/core/pirtm\_link\_ensemble.py}: linkage between PIRTM recurrences and ensemble states.
  \item \texttt{meta\_ensembles/tests}: a suite of tests validating ledger behavior, aggregation, and PIRTM-linked folds.
  \item \texttt{meta\_ensembles/cli.py}: an MVP CLI surface.
\end{itemize}

Originally, the test suite failed due to unresolved \texttt{pirtm} imports and fixture path mismatches. The goal was to restore a green baseline without collapsing the monorepo structure.

\subsection{Pytest Configuration for Local PIRTM}

A new \texttt{conftest.py} in \texttt{packages/meta-ensembles} ensures that the local \texttt{mirror-dissonance-pro/pirtm} tree is importable when running tests:

\begin{lstlisting}[style=code,language=Python]
# packages/meta-ensembles/conftest.py

from __future__ import annotations

import sys
from pathlib import Path

ROOT = Path(__file__).resolve().parents[2]
PIRTM_ROOT = ROOT / "mirror-dissonance-pro" / "pirtm"

# Insert the parent of the pirtm package so "import pirtm" works.
sys.path.insert(0, str(PIRTM_ROOT.parent))
\end{lstlisting}

This approach:
\begin{itemize}
  \item Keeps sys.path manipulation local to tests, not package code.
  \item Matches patterns already in use for other HQ packages (e.g., \texttt{sigma-kernel}).
  \item Allows \texttt{import pirtm} to behave as expected without requiring external installation.
\end{itemize}

\subsection{Fixture Path Correction}

The test \texttt{meta\_ensembles/tests/test\_meta\_fold.py} referenced PIRTM fixtures from an incorrect path. It has been updated to resolve fixtures relative to the HQ repo root, where PIRTM's fixtures live under \texttt{mirror-dissonance-pro/pirtm/tests/fixtures}. Structurally, the change looks like:

\begin{lstlisting}[style=code,language=Python]
from pathlib import Path

ROOT = Path(__file__).resolve().parents[3]
FIXTURES_DIR = ROOT / "mirror-dissonance-pro" / "pirtm" / "tests" / "fixtures"

def load_pirtm_fixture(name: str) -> str:
    return (FIXTURES_DIR / name).read_text()
\end{lstlisting}

With this change, running:

\begin{lstlisting}[style=code]
cd packages/meta-ensembles
python3 -m pytest -q
\end{lstlisting}

reports \emph{81 passed}, indicating a clean MVP baseline for further meta-ensemble development.

\subsection{Ensemble Ledger and PIRTM Link: Conceptual Sketch}

While the actual implementation details are in the codebase, the conceptual structure can be captured as follows.

Let \(\mathcal{M}\) be a set of models or oracles, and let \(x\) denote an input (e.g., a governance scenario). Each model \(m \in \mathcal{M}\) outputs a judgment \(y_m \in \mathcal{Y}\). The ensemble ledger is a sequence of records:
\[
  L = \{(x_i, \{y_{m,i}\}_{m\in\mathcal{M}}, t_i, \text{metadata}_i)\}_{i=1}^N.
\]

The meta-ensemble aggregator computes an ensemble decision:
\[
  y^{\star}_i = A\left(\{y_{m,i}\}_{m\in\mathcal{M}}, \text{weights}, \text{constraints}\right).
\]

A PIRTM link ensemble then interprets the ensemble states (e.g., probability distributions, tensor encodings of decisions) as part of a recursive tensor dynamical system:
\[
  s_{k+1} = F(s_k; \theta, p_k)
\]
where:
\begin{itemize}
  \item \(s_k\) is the ensemble state at prime index \(p_k\),
  \item \(\theta\) are parameters of the meta-ensemble dynamics,
  \item \(F\) is defined via PIRTM's contractive recurrence runtime.
\end{itemize}

In code, a simplified PIRTM-linked ensemble iteration could look like:

\begin{lstlisting}[style=code,language=Python]
import numpy as np
from meta_ensembles.core.ensemble_ledger import EnsembleLedger
from meta_ensembles.core.pirtm_link_ensemble import PirtmLinkedEnsemble

ledger = EnsembleLedger()
linked = PirtmLinkedEnsemble(ledger=ledger)

state = linked.initial_state()
for p_idx, record in enumerate(ledger.iter_prime_indexed()):
    state = linked.step(state, record, prime_index=p_idx)

final_state = linked.project_decision(state)
\end{lstlisting}

A full mathematical treatment would express the prime-indexing and contraction properties explicitly as part of a multiplicity-theoretic framework, but this report focuses on the integration scaffolding.

\section{Governance and Open-Core Considerations}

\subsection{L0 Invariants and Optional Dependencies}

Phase Mirror's governance documents emphasize that L0 invariants are:
\begin{itemize}
  \item Always-on, non-configurable.
  \item Cheap enough that not running them costs more in missed bugs than running them.
  \item Independent of external research or experimental features \cite{pm-mvp-analysis,pm-open-core-dev}.
\end{itemize}

PIRTM integration, meta-ensembles, and future MLIR-based optimizations must therefore:
\begin{itemize}
  \item Never be required for L0 invariant evaluation.
  \item Be clearly marked experimental/optional in ADRs and user-facing docs.
  \item Provide graceful degradation paths and explicit failure signaling (degraded mode) rather than implicit breakage.
\end{itemize}

\subsection{Open-Core Boundary and Research Velocity}

Phase Mirror's open-core model reserves certain features---Tier B semantic rules, FP calibration networks, compliance packs---for a pro tier, but ensures the open-core remains genuinely useful on its own \cite{pm-open-core-dev,pm-mvp-analysis}. PIRTM and meta-ensembles fit into this picture as:
\begin{itemize}
  \item Mathematical and algorithmic engines that can live in either open-core or pro domains depending on how they are applied.
  \item Potential sources of research velocity (new rules, better dissonance metrics) if integrated with calibration stores and ensemble statistics.
\end{itemize}

The key is to avoid entangling open-core adoption with proprietary configuration or pro-only dependencies. Moving \texttt{pirtm} to an external, properly versioned package, and ensuring \texttt{meta-ensembles} can run without proprietary trees, are concrete steps in this direction.

\section{Recommended Next Steps and Precision Questions}

\subsection{Levers and Metrics}

We summarize a few concrete levers for the immediate future:

\begin{itemize}
  \item \textbf{Runtime-first PIRTM integration}
    \begin{itemize}
      \item Owner: Phase Mirror Python MVP maintainer.
      \item Metric: All \texttt{phase-mirror-mvp} tests pass with \texttt{pirtm} installed from an external source, no \texttt{sys.path} hacks.
      \item Horizon: 7--10 days.
    \end{itemize}
  \item \textbf{MLIR optional backend}
    \begin{itemize}
      \item Owner: PIRTM maintainer.
      \item Metric: Separate CI job demonstrating MLIR-backed runs; Phase Mirror still passes tests when MLIR is absent.
      \item Horizon: 30 days.
    \end{itemize}
  \item \textbf{Meta-ensembles decoupling}
    \begin{itemize}
      \item Owner: Meta-ensembles maintainer.
      \item Metric: \texttt{meta-ensembles} tests can run in a mode where PIRTM is mocked or optional, and a mode where real PIRTM is used.
      \item Horizon: 14 days.
    \end{itemize}
\end{itemize}

\subsection{Precision Questions}

Before locking in the long-term integration pattern, we recommend explicitly resolving:

\begin{enumerate}[label=(Q\arabic*)]
  \item Are we willing to codify as an L0-style invariant that \emph{Phase Mirror governance decisions must never depend exclusively on MLIR availability}? This would permanently bound the blast radius of MLIR-level failures.
  \item In terms of release sequencing, should PIRTM-backed tensor recurrences be part of the success criteria for the next Phase Mirror public milestone, or remain explicitly non-blocking to preserve MVP scope?
  \item For meta-ensembles, is the primary goal near-term \emph{diagnostic diversity} (ensembles as multiple mirrors on dissonance) or \emph{performance/accuracy} (ensembles as better predictors)? This affects how we prioritize ledger semantics vs. tensor dynamics.
\end{enumerate}

\section{Conclusion}

The work described in this report brings the PIRTM, Phase Mirror, and meta-ensembles projects into a tighter, but \emph{explicitly governed}, relationship. By:
\begin{itemize}
  \item Implementing a runtime-first PIRTM adapter and corresponding tests,
  \item Formalizing integration decisions via ADR-028,
  \item Stabilizing the meta-ensembles package with a clean \texttt{81 passed} test baseline,
\end{itemize}
the system is now well-positioned to:
\begin{itemize}
  \item Transition to an external \texttt{pirtm} package boundary,
  \item Introduce MLIR as an optional backend with clear failure semantics,
  \item Explore multiplicity-theoretic ensembles within a governance-safe envelope.
\end{itemize}

The central challenge, and opportunity, is to keep the mathematical richness of PIRTM and meta-ensembles in tension with Phase Mirror's governance invariants, using ADRs, adapters, and CI policies as the binding artifacts.

\appendix
\section*{Appendix: Contractive Recurrences and Operator Norm Bounds}

\section{Abstract Contractive PIRTM Recurrence}

We model the PIRTM runtime used in the Phase Mirror integration as a discrete-time recurrence on a Banach space, abstracting away from implementation details (NumPy arrays, concrete update rules) while capturing the essential contractive structure that guarantees stability of the dynamics.[pirtm-pyproject][hq-pirtm-local]

Let \((X,\|\cdot\|)\) be a real Banach space, typically \(X = \mathbb{R}^d\) with the Euclidean norm induced by NumPy tensors.[pirtm-pyproject] Let \(\Theta\) be a parameter space (e.g., a finite-dimensional vector space of coefficients).

\begin{definition}[PIRTM-style contractive operator]
A mapping \(F : X \times \Theta \to X\) is said to be \emph{globally contractive} in its first argument with constant \(0 \leq \lambda < 1\) if
\[
  \|F(x,\theta) - F(y,\theta)\| \leq \lambda \, \|x - y\|
\]
for all \(x,y \in X\) and all \(\theta \in \Theta\).[hq-pirtm-local]
\end{definition}

In the PIRTM recurrence runtime, a typical iteration has the form
\[
  x_{k+1} = F(x_k,\theta), \quad k = 0,1,2,\dots,
\]
where \(\theta\) may encode prime-indexed structure, contractive weights, or other multiplicity-theoretic features, but for the purposes of stability we assume only that \(F(\cdot,\theta)\) is contractive.[hq-pirtm-local]

\begin{theorem}[Existence and uniqueness of fixed point; convergence rate]
\label{thm:banach}
Suppose \(F(\cdot,\theta)\) is globally contractive with constant \(\lambda \in [0,1)\). Then:
\begin{enumerate}
  \item There exists a unique fixed point \(x^\star \in X\) such that \(F(x^\star,\theta) = x^\star\).
  \item For any initial state \(x_0 \in X\), the sequence defined by \(x_{k+1} = F(x_k,\theta)\) converges to \(x^\star\).
  \item The convergence is \emph{geometric}: for all \(k \geq 0\),
  \[
    \|x_k - x^\star\| \leq \lambda^k \, \|x_0 - x^\star\|.
  \]
\end{enumerate}
\end{theorem}

\begin{proof}
Define \(T := F(\cdot,\theta):X\to X\). By assumption,
\[
  \|T(x) - T(y)\| \leq \lambda \|x-y\|,
\]
so \(T\) is a contraction mapping on the complete metric space \((X,\|\cdot\|)\). By the Banach fixed point theorem, there exists a unique \(x^\star \in X\) such that \(T(x^\star) = x^\star\), and for any \(x_0 \in X\), the iterates \(x_k = T^k(x_0)\) converge to \(x^\star\).[hq-pirtm-local]

For the geometric bound, we have for all \(k\geq 0\),
\[
  \|x_{k+1} - x^\star\|
  = \|T(x_k) - T(x^\star)\|
  \leq \lambda \|x_k - x^\star\|.
\]
By induction, \(\|x_k - x^\star\| \leq \lambda^k \|x_0 - x^\star\|\), which completes the proof.
\end{proof}

\subsection{Operator Norm Interpretation}

In the linear case, when \(F(x,\theta) = A(\theta)x + b(\theta)\) with a bounded linear operator \(A(\theta)\) and bias \(b(\theta)\), the contractive condition can be stated in terms of the operator norm:
\[
  \|A(\theta)\|_{\text{op}} := \sup_{x \neq 0} \frac{\|A(\theta)x\|}{\|x\|} < 1.
\]

\begin{proposition}[Contractivity via operator norm]
If \(A(\theta)\) is linear and \(\|A(\theta)\|_{\text{op}} \leq \lambda < 1\), then \(F(x,\theta) = A(\theta)x + b(\theta)\) is globally contractive with constant \(\lambda\). Consequently, the PIRTM recurrence
\[
  x_{k+1} = A(\theta)x_k + b(\theta)
\]
has a unique fixed point \(x^\star\), and the geometric convergence bound in Theorem~\ref{thm:banach} holds.
\end{proposition}

\begin{proof}
For any \(x,y \in X\),
\[
  \|F(x,\theta) - F(y,\theta)\|
  = \|A(\theta)x + b(\theta) - (A(\theta)y + b(\theta))\|
  = \|A(\theta)(x-y)\|
  \leq \|A(\theta)\|_{\text{op}} \,\|x-y\|
  \leq \lambda \|x-y\|.
\]
Thus \(F(\cdot,\theta)\) is a contraction with constant \(\lambda\) and Theorem~\ref{thm:banach} applies.
\end{proof}

\section{Prime-Indexed Recurrence and Norm Bounds}

PIRTM emphasises prime-indexed recursions that encode multiplicity across non-linear scales.[hq-pirtm-local] Let \((p_k)_{k\geq 1}\) denote the increasing sequence of prime numbers and consider a state sequence \((x_{p_k})_{k\geq 1}\) defined by
\[
  x_{p_{k+1}} = F_{p_{k+1}}(x_{p_k}), \quad k \geq 1,
\]
where each \(F_{p_k}:X\to X\) is a contraction with constant \(\lambda_{p_k} \in [0,1)\).

\begin{assumption}[Uniform prime-indexed contractivity]
There exists a constant \(\lambda^\star < 1\) such that
\[
  \sup_{k \geq 1} \lambda_{p_k} \leq \lambda^\star.
\]
\end{assumption}

Under this assumption, we obtain a uniform bound on the deviation between two prime-indexed trajectories.

\begin{proposition}[Uniform stability across prime indices]
Let \((x_{p_k})_{k\geq 1}\) and \((y_{p_k})_{k\geq 1}\) be sequences defined by
\[
  x_{p_{k+1}} = F_{p_{k+1}}(x_{p_k}), \quad
  y_{p_{k+1}} = F_{p_{k+1}}(y_{p_k}),
\]
with the same family \(\{F_{p_k}\}\) and initial states \(x_{p_1},y_{p_1} \in X\). If the uniform contractivity assumption holds, then for all \(k\geq 1\),
\[
  \|x_{p_k} - y_{p_k}\| \leq (\lambda^\star)^{k-1} \, \|x_{p_1} - y_{p_1}\|.
\]
\end{proposition}

\begin{proof}
For each step,
\[
  \|x_{p_{k+1}} - y_{p_{k+1}}\|
  = \|F_{p_{k+1}}(x_{p_k}) - F_{p_{k+1}}(y_{p_k})\|
  \leq \lambda_{p_{k+1}} \|x_{p_k} - y_{p_k}\|
  \leq \lambda^\star \|x_{p_k} - y_{p_k}\|.
\]
Iterating from \(k=1\) yields
\[
  \|x_{p_k} - y_{p_k}\| \leq (\lambda^\star)^{k-1} \|x_{p_1} - y_{p_1}\|.
\]
\end{proof}

This inequality gives a multiplicity-theoretic version of contractive stability: even when the update operator varies with the prime index, a uniform operator norm bound ensures geometric decay of differences.

\section{Meta-Ensembles as Contractive Lifts}

In the meta-ensembles package, ensemble state updates can be seen as lifting base-model decisions into a higher-dimensional recurrence.[hq-pirtm-local] Abstractly, suppose:
\begin{itemize}
  \item \(\mathcal{M} = \{1,\dots,M\}\) indexes base models or oracles.
  \item For each \(i \in \mathcal{M}\), \(f_i : X \to Y\) is a decision map (e.g., \(Y\) is a finite label set or a probability simplex).
  \item The ensemble state space is \(Z = \mathbb{R}^M\) or \(Z = \mathbb{R}^{M \times d}\), embedding model outputs.
\end{itemize}

Define an ensemble embedding \(E : Y^\mathcal{M} \to Z\) and an aggregation map \(A : Z \to Y\). A meta-ensemble iteration coupled to PIRTM can be written as:
\[
  z_{k+1} = G(z_k; \theta), \quad
  y_k = A(z_k),
\]
where \(G : Z \to Z\) is itself defined by a PIRTM-style contractive operator at the ensemble level.

\begin{assumption}[Ensemble-level contractivity]
There exists \(0 \leq \mu < 1\) such that
\[
  \|G(z) - G(w)\| \leq \mu \, \|z - w\|
\]
for all \(z,w \in Z\), with \(\|\cdot\|\) an appropriate norm (e.g., Euclidean or maximum norm).
\end{assumption}

\begin{proposition}[Stability of meta-ensemble iterates]
Under ensemble-level contractivity, the meta-ensemble state sequence \((z_k)\) admits a unique fixed point \(z^\star\), and for all \(k\geq 0\),
\[
  \|z_k - z^\star\| \leq \mu^k \, \|z_0 - z^\star\|.
\]
If the aggregation map \(A:Z\to Y\) is Lipschitz with constant \(L_A\), then the induced decisions \(y_k = A(z_k)\) satisfy
\[
  \|y_k - y^\star\| \leq L_A \mu^k \, \|z_0 - z^\star\|,
\]
where \(y^\star = A(z^\star)\).
\end{proposition}

\begin{proof}
The first inequality follows by applying Theorem~\ref{thm:banach} with \(X:=Z\) and \(F:=G\). For the decisions, we have
\[
  \|y_k - y^\star\| = \|A(z_k) - A(z^\star)\| \leq L_A \|z_k - z^\star\|
  \leq L_A \mu^k \|z_0 - z^\star\|.
\]
\end{proof}

In the concrete \texttt{meta-ensembles} implementation, the ledger and PIRTM link ensemble modules respect this pattern: ensemble state updates are implemented via well-behaved transforms (typically finite-dimensional linear or Lipschitz non-linear maps), and tests verify stability and boundedness of the resulting iterates.[hq-pirtm-local]

\section{Error Propagation and Norm Bounds for Adapters}

Phase Mirror’s adapter layer (e.g., FP store, secret store, consent store) adopts a ``throw, don’t mask’’ error contract.[file:29][file:40] This can also be expressed in operator-norm style bounds for the \emph{perturbed} recurrence driven by occasional adapter failures.

Consider a recurrence
\[
  x_{k+1} = F(x_k) + e_k,
\]
where \(F : X \to X\) is contractive with constant \(\lambda < 1\) and \((e_k)_{k\geq 0} \subset X\) captures perturbations from adapter errors that are caught and mapped into explicit, bounded deviations before being handled at L0.[file:29] Assume
\[
  \|e_k\| \leq \varepsilon_k
\]
for some non-negative sequence \((\varepsilon_k)\).

\begin{proposition}[Stability under bounded perturbations]
\label{prop:perturbed}
Let \(x^\star\) be the fixed point of \(F\). Then the perturbed recurrence satisfies
\[
  \|x_{k+1} - x^\star\| \leq \lambda \|x_k - x^\star\| + \varepsilon_k
\]
for all \(k\geq 0\), and hence
\[
  \|x_k - x^\star\|
  \leq \lambda^k \|x_0 - x^\star\| + \sum_{j=0}^{k-1} \lambda^{k-1-j} \varepsilon_j.
\]
\end{proposition}

\begin{proof}
We have
\begin{align*}
  \|x_{k+1} - x^\star\|
  &= \|F(x_k) + e_k - F(x^\star)\| \\
  &\leq \|F(x_k) - F(x^\star)\| + \|e_k\| \\
  &\leq \lambda \|x_k - x^\star\| + \varepsilon_k.
\end{align*}
Unrolling this inequality yields
\[
  \|x_k - x^\star\|
  \leq \lambda^k \|x_0 - x^\star\| + \sum_{j=0}^{k-1} \lambda^{k-1-j} \varepsilon_j
\]
by induction on \(k\).
\end{proof}

\begin{corollary}[Uniform perturbation bound]
If \(\varepsilon_k \leq \overline{\varepsilon}\) for all \(k\), then
\[
  \|x_k - x^\star\|
  \leq \lambda^k \|x_0 - x^\star\|
       + \overline{\varepsilon} \frac{1-\lambda^k}{1-\lambda}.
\]
In particular,
\[
  \limsup_{k\to\infty} \|x_k - x^\star\|
  \leq \frac{\overline{\varepsilon}}{1-\lambda}.
\]
\end{corollary}

This provides a quantitative version of Phase Mirror’s ``fail-closed, observable’’ philosophy at the recurrence level: as long as the magnitude of adapter-induced perturbations remains bounded (and relatively small compared to \(1-\lambda\)), the system remains close to the ideal fixed point and deviations are controllable.[file:29]

\section{Summary of Bounds}

For convenience, we collect the key inequalities derived in this appendix.

\begin{itemize}
  \item \textbf{Pure contraction, fixed operator:}
  \[
    \|x_k - x^\star\| \leq \lambda^k \|x_0 - x^\star\|.
  \]
  \item \textbf{Linear case via operator norm:}
  \[
    \|A\|_{\mathrm{op}} < 1 \;\Rightarrow\; \text{contraction with constant }\lambda=\|A\|_{\mathrm{op}}.
  \]
  \item \textbf{Prime-indexed varying operators (uniformly contractive):}
  \[
    \|x_{p_k} - y_{p_k}\| \leq (\lambda^\star)^{k-1} \|x_{p_1} - y_{p_1}\|.
  \]
  \item \textbf{Meta-ensemble lift (ensemble-level contraction):}
  \[
    \|z_k - z^\star\| \leq \mu^k \|z_0 - z^\star\|,\quad
    \|y_k - y^\star\| \leq L_A \mu^k \|z_0 - z^\star\|.
  \]
  \item \textbf{Perturbed recurrence with bounded noise:}
  \[
    \|x_k - x^\star\|
    \leq \lambda^k \|x_0 - x^\star\|
       + \sum_{j=0}^{k-1} \lambda^{k-1-j} \varepsilon_j,
  \]
  and under uniform perturbations,
  \[
    \limsup_{k\to\infty} \|x_k - x^\star\|
    \leq \frac{\overline{\varepsilon}}{1-\lambda}.
  \]
\end{itemize}

These bounds are the mathematical backbone behind the informal ``contractive recurrence runtime’’ phrasing for PIRTM and the stability guarantees implicitly relied upon in the Phase Mirror and meta-ensembles integrations.[pirtm-pyproject][hq-pirtm-local][file:29][file:40]

\nocite{*}
\bibliographystyle{unsrt}  
\bibliography{references}
\end{document}